\chapter{Sociétés de trading pour compte propre à moyenne fréquence}\label{ch:in:medium-frequency-proprietary-shops}

La plupart des sociétés de trading pour compte propre sont petites. Fin 2024, le régulateur américain des broker-dealers a classé 103 des 3\,249
broker-dealers qu'il supervise dans ses segments d'activité de trading pour compte propre et de market making~; 93 d'entre eux comptaient 150
représentants inscrits ou moins, et aucun n'en comptait 500. Parmi l'ensemble des broker-dealers, près de la moitié comptait dix personnes inscrites
ou moins, tandis que les 149 plus grands employaient 82~\% de l'ensemble des représentants inscrits. Les sociétés de mille personnes des
chapitres~2 à~4 forment la queue d'une longue distribution dont le corps regroupe des sociétés de quelques dizaines de personnes, beaucoup d'entre elles opérant
à des horizons de quelques minutes à quelques jours. Ce chapitre décrit ce corps~: ce que signifie la moyenne fréquence, comment les tailles des sociétés
se répartissent, comment les sociétés entrent et sortent, et les arrangements --- partages de profits, \omterm{def:in:medium-frequency-proprietary-shops:arcade}{salles de trading partagées}, \omterm{def:in:medium-frequency-proprietary-shops:funded}{programmes de trader financé} --- que
rencontre un candidat à l'emploi du côté des petites structures, dont certains ne sont pas du tout des emplois.

\section{Moyenne fréquence~: horizons, capacité et infrastructure}

\begin{definition}[Trading à moyenne fréquence]\label{def:in:medium-frequency-proprietary-shops:mft}
Le \emph{trading à moyenne fréquence}\index{trading à moyenne fréquence} est un trading systématique dont les positions durent de quelques minutes à
quelques jours~: plus lent que le trading à haute fréquence, dont l'edge se mesure en microsecondes et les durées de détention en secondes
(livre~10, chapitre~24), et plus rapide que les signaux quotidiens à mensuels de la plupart des fonds systématiques.
\end{definition}

L'horizon détermine la société. Une stratégie à moyenne fréquence rapporte davantage par transaction qu'un teneur de marché et négocie moins
souvent, si bien que ses coûts sont l'impact de marché et les frais plutôt que le prix d'être le premier~; son infrastructure exige de bonnes données
et une bonne exécution (livre~10), mais pas l'ingénierie tick-to-trade du livre~13~; et sa capacité, c'est-à-dire le capital qu'elle peut gérer avant
que ses propres opérations n'érodent son edge (capacité de stratégie, livre~7, chapitre~28), est supérieure à celle d'un teneur de marché et inférieure à
celle d'un trend follower. Sa recherche est celle du livre~7 et ses stratégies celles du livre~8~: prédicteurs intrajournaliers et multi-jours sur actions,
sur contrats à terme et devises. Une petite équipe disposant de bonnes données et d'un processus de recherche rigoureux peut faire tourner un tel portefeuille, ce qui explique
pourquoi la catégorie compte tant de petites sociétés.

\begin{remark}[Pour compte propre, pas nécessairement petit]\label{rem:in:medium-frequency-proprietary-shops:large}
Les grandes sociétés opèrent aussi à moyenne fréquence, souvent au sein de fonds multistratégies et de plateformes (chapitres~4 et~5). Le
présent chapitre porte sur les sociétés de compte propre indépendantes, pour lesquelles c'est toute l'activité~: des sociétés qui ne négocient que leurs propres
capitaux et ceux de leurs partenaires.
\end{remark}

\section{La distribution des tailles des sociétés de compte propre}

Le seul recensement public des tailles de sociétés de trading est celui du régulateur des broker-dealers, et son unité est le
représentant inscrit~: une personne inscrite pour traiter avec le public, et non un salarié. Une société de compte propre sans clients peut
inscrire peu de son personnel. Le recensement dit néanmoins comment sont structurées les sociétés du secteur.

\begin{dated}{2024-12}{Broker-dealers par taille}\label{dat:in:medium-frequency-proprietary-shops:sizes}
FINRA, fin 2024~: 3\,249 broker-dealers inscrits. Selon le nombre de représentants inscrits~: 1\,522 sociétés comptant dix ou moins,
328 comptant 11--15, 205 comptant 16--20, 134 comptant 21--25, 108 comptant 26--30, 159 comptant 31--40, 102 comptant 41--50, 143 comptant 51--75, 80 comptant
76--100, 110 comptant 101--150, 137 comptant 151--300, 72 comptant 301--499, 65 comptant 500--1\,000 et 84 comptant plus de 1\,000. Sur les
649\,241 représentants inscrits, 81.7~\% travaillent dans les 149 grandes sociétés (500 ou plus), 8.9~\% dans les sociétés de taille moyenne et
9.4~\% dans les petites. Dans ses segments d'activité, la FINRA compte 93 petites sociétés et 10 sociétés de taille moyenne
dans le trading pour compte propre et le market making, et aucune parmi les grandes.
\end{dated}

La plupart des sociétés sont petites et la plupart des personnes travaillent dans les grandes~: 46.8~\% des broker-dealers comptent dix représentants inscrits
ou moins, et 78.7~\% en comptent cinquante ou moins, mais ces petites sociétés regroupent 9.4~\% des représentants.
Une distribution de cette forme se décrit en général par sa queue supérieure.

\begin{method}[Ajustement d'une queue sur des effectifs regroupés]\label{met:in:medium-frequency-proprietary-shops:fit}
\begin{enumerate}
\item Choisir une borne inférieure $x_{\min}$ et garder les classes supérieures ou égales à celle-ci.
\item Écrire la probabilité de chaque classe sous une queue de Pareto $P(X>x)=(x/x_{\min})^{-\alpha}$, en traitant les bornes entières comme
continues aux demi-unités près~; la log-vraisemblance vaut $\sum_i n_i\ln p_i(\alpha)$.
\item Maximiser par rapport à $\alpha$~; l'erreur type est l'inverse de la racine carrée de l'opposé de la dérivée seconde de la
log-vraisemblance au maximum.
\item Comparer à une autre forme (une lognormale) par la distance du khi-deux entre effectifs observés et théoriques, et
déplacer $x_{\min}$ pour voir si $\alpha$ est stable.
\end{enumerate}
\end{method}

À partir de onze représentants inscrits et au-delà (1\,727 sociétés, treize classes), l'indice de queue ajusté vaut $\alpha=0.60$ avec une
erreur type de 0.015 (la \cref{fig:in:medium-frequency-proprietary-shops:ccdf}). Un indice inférieur à un signale une queue lourde~: dans
une telle distribution, les plus grandes sociétés détiennent une part du total qui ne diminue pas à mesure qu'on ajoute des sociétés, ce que
la répartition des représentants l'indique directement. L'ajustement est imparfait. La droite de Pareto surestime la part des
très grandes sociétés (6.5~\% des sociétés au-delà de 1\,000 représentants, contre 4.9~\% observées), une lognormale ajuste mieux les classes (une
distance du khi-deux de 5.1 contre 21.2), et l'indice passe de 0.60 à 0.65 à mesure que la borne inférieure passe de 11 à 51.
La lecture honnête est une fourchette, non un nombre~: un indice de queue compris entre environ 0.6 et 0.65, soit une queue lourde quelle que soit la valeur retenue.

\begin{omfigure}
\begin{tikzpicture}
\begin{axis}[width=10.5cm,height=6cm,xmode=log,ymode=log,xmin=9,xmax=1500,ymin=0.03,ymax=1.2,log ticks with fixed point,
  xlabel={\small représentants inscrits, $x$ (échelle log)},ylabel={\small part des sociétés ayant au moins $x$ (échelle log)},
  tick label style={font=\footnotesize},grid=major,grid style={gray!20},table/col sep=comma,
  legend style={legend cell align=left,font=\scriptsize,at={(0.5,-0.3)},anchor=north,/tikz/every even column/.append style={column sep=8pt}},legend columns=2]
\addplot[only marks,mark=*,omDef] table[x=size,y=observed]{figdata/industry/06-medium-frequency-proprietary-shops/ccdf.csv};
\addplot[omThm,thick] table[x=size,y=pareto]{figdata/industry/06-medium-frequency-proprietary-shops/ccdf.csv};
\legend{observées (classes FINRA),{queue de Pareto, $\alpha=0.60$}}
\end{axis}
\end{tikzpicture}

\omcaption{Part des 1\,727 broker-dealers comptant onze représentants inscrits ou plus ayant au moins $x$,
fin 2024, en échelles logarithmiques, avec la queue de Pareto ajustée. Données~: FINRA Industry Snapshot 2025, tableau 2.1.9, via
\texttt{in\_firmsize.fit}.}\label{fig:in:medium-frequency-proprietary-shops:ccdf}
\end{omfigure}

\section{Nouveaux entrants et sociétés issues de la crypto}

Les sociétés entrent et sortent de la population inscrite chaque année
(la \cref{fig:in:medium-frequency-proprietary-shops:flows}). Entre 2011 et 2024, en moyenne 5.6~\% des broker-dealers de l'année précédente
ont cessé d'être membres de la FINRA et 3.2~\% y sont entrés~; en 2024, 185 ont quitté et 136 sont entrés.
La population a diminué de 29~\% de 2010 à 2024 parce que les sorties ont dépassé les entrées chaque année, et non parce que les entrées se sont arrêtées. Parmi
les sociétés de compte propre, un nouvel entrant peut être une poignée de traders et de chercheurs disposant de leur propre capital ou de celui d'un bailleur de fonds,
et une sortie peut être une société qui ferme, fusionne ou quitte le périmètre du régulateur en ne négociant que des contrats à terme ou que des actifs
crypto, activité que l'inscription comme broker-dealer ne couvre pas.

Les sociétés issues des marchés crypto sont le cas inverse de ce dernier. Les sociétés de trading crypto du chapitre~10 ont commencé
en dehors des registres des régulateurs de valeurs mobilières, en trading vingt-quatre heures sur vingt-quatre sur des plateformes de négociation dotées de leurs propres règles (livre~3, chapitre~24),
et une société crypto qui s'étend aux contrats à terme,
aux actions ou aux options doit obtenir les inscriptions que cela exige. Pour un salarié
leur attrait est celui de toute jeune société~: une petite équipe, de larges responsabilités et une rémunération qui dépend fortement des
résultats propres de la société.

\begin{omfigure}
\begin{tikzpicture}
\begin{axis}[width=11cm,height=5.2cm,xmin=2010.5,xmax=2024.5,ymin=0,ymax=8,ylabel={\small \% des sociétés de l'année précédente},
  xtick={2011,2014,2017,2020,2023},xticklabel style={/pgf/number format/1000 sep={}},tick label style={font=\footnotesize},
  grid=major,grid style={gray!20},table/col sep=comma,
  legend style={legend cell align=left,font=\scriptsize,at={(0.5,-0.22)},anchor=north,/tikz/every even column/.append style={column sep=8pt}},legend columns=2]
\addplot[omThm,thick,mark=*] table[x=year,y=exit_pct]{figdata/industry/06-medium-frequency-proprietary-shops/flows.csv};
\addplot[omDef,thick,mark=square*] table[x=year,y=entry_pct]{figdata/industry/06-medium-frequency-proprietary-shops/flows.csv};
\legend{sociétés qui sortent,sociétés qui entrent}
\end{axis}
\end{tikzpicture}

\omcaption{Broker-dealers qui quittent et rejoignent chaque année l'affiliation à la FINRA, 2011--2024, en pourcentage de la population
de fin d'année précédente. Données~: FINRA Industry Snapshot 2025, tableau 2.2.3, via
\texttt{in\_firmsize.entry\_exit}.}\label{fig:in:medium-frequency-proprietary-shops:flows}
\end{omfigure}

\section{Partage des profits, salles de trading partagées et programmes de trader financé}

Au bas de l'échelle, la frontière entre un emploi et une activité à son compte devient floue. Trois dispositifs reviennent régulièrement.

\textbf{Le partage des profits.} Un trader d'une petite société de trading pour compte propre est souvent rémunéré par une part du P\&L qu'il génère, net de ses coûts,
avec peu ou pas de salaire~: une rémunération à la formule au sens du livre~16 (chapitre~10), la société apportant en contrepartie son capital,
sa technologie et sa gestion des risques, qui constituent la contrepartie de la société. La part, les coûts facturés et le report ou non des pertes
déterminent ce que gagne le trader (chapitre~13).

\begin{definition}[Salle de trading partagée (trading arcade)]\label{def:in:medium-frequency-proprietary-shops:arcade}
Une \emph{salle de trading partagée}\index{salle de trading partagée} (trading arcade) est une structure dans laquelle plusieurs entités juridiques ou personnes physiques négocient
sous la licence d'une seule société pour le trading pour compte propre, en général chacune avec sa propre allocation de capital,
tout en partageant l'infrastructure, la gestion des risques et la conformité du titulaire de la licence.
\end{definition}

Le régulateur néerlandais décrit les arcades en ces termes et indique qu'elles sont possibles en principe, à condition que toutes
les exigences applicables aux entreprises d'investissement soient respectées, les droits, pouvoirs et obligations de chaque participant étant clairement
attribués et documentés. Les règles de fonds propres qui s'imposent au titulaire de la licence (fonds propres et exigence au titre des facteurs K, livre~16,
chapitre~2) s'appliquent à l'ensemble de l'arcade~: les pertes d'un participant deviennent le problème du titulaire de la licence.

\begin{definition}[Programme de trader financé (funded-trader programme)]\label{def:in:medium-frequency-proprietary-shops:funded}
Un \emph{programme de trader financé}\index{programme de trader financé} (funded-trader programme) est une activité qui vend à des clients particuliers la possibilité de
négocier le capital de l'opérateur~: les clients paient des frais pour passer une évaluation, en général sur un compte simulé, et ceux qui
atteignent ses objectifs se voient promettre une part des profits d'un compte «~financé~».
\end{definition}

Un \omterm{def:in:medium-frequency-proprietary-shops:funded}{programme de trader financé} n'est pas un employeur, et ses clients ne sont pas des traders d'une société de trading pour compte propre~: ils versent
des frais à l'opérateur, qui tire son revenu de ces frais. Le dossier public montre comment une telle activité peut se retrouver
devant un régulateur, et comment une telle affaire peut se terminer.

\begin{dated}{2025-05}{Une affaire de trader financé dans le dossier judiciaire}\label{dat:in:medium-frequency-proprietary-shops:mff}
En août 2023, la CFTC a attaqué l'opérateur d'un \omterm{def:in:medium-frequency-proprietary-shops:funded}{programme de trader financé} devant le tribunal de district des États-Unis pour le district du New
Jersey, en affirmant que plus de 135\,000 clients s'étaient inscrits, en versant au moins 310~millions de dollars de frais~; que certains
types de comptes exigeaient d'abord que les clients réalisent des profits simulés sur un compte de démonstration~; et que les comptes financés
étaient présentés avec la promesse de «~jusqu'à 85~\% des profits~». Le 13 mai 2025, un special master désigné par le tribunal, statuant sur la
requête des défendeurs en sanctions contre la CFTC, a recommandé d'accueillir la requête et de rejeter la plainte
avec préjudice. Le tribunal a rejeté l'affaire avec préjudice et a accordé aux défendeurs les frais et dépens de la
requête en sanctions. Les allégations n'ont jamais été prouvées.
\end{dated}

\section{Ce que signifie une petite société pour le salarié}

Un emploi dans une petite société de trading pour compte propre diffère de celui d'une grande société sur quatre points qui découlent des chiffres présentés plus haut.
\begin{enumerate}
\item \textbf{Le capital appartient aux associés.} Les pertes réduisent ce que possèdent les propriétaires, et l'exigence de capital de la société
(livre~16, chapitre~2) limite le risque que l'ensemble de la société peut prendre~; un mauvais mois peut réduire l'allocation de chacun.
\item \textbf{La rémunération suit de près le P\&L.} Partages à la formule, peu de postes salariés et une rémunération différée souvent absente.
\item \textbf{Les contrôles sont plus légers.} Risques, conformité et opérations peuvent tenir chacun à une seule personne~; le salarié en fait davantage de
tout, y compris la vérification de son propre travail.
\item \textbf{La société peut ne pas durer.} Environ un broker-dealer sur dix-huit quitte le registre chaque année. La valeur d'option de l'emploi tient à ce qu'il enseigne et aux relations qu'il ouvre.
\end{enumerate}

\begin{example}[Le partage d'un trader]\label{ex:in:medium-frequency-proprietary-shops:split}
Un trader d'une petite société conserve 40~\% du P\&L net de son portefeuille après des coûts de 150\,000~\$ par an, sans salaire. Une année de
800\,000~\$ de P\&L brut lui rapporte $0.4\times(800\,000-150\,000)=\$260\,000$~; une année de 200\,000~\$ lui rapporte 20\,000~\$~; une année de
100\,000~\$ ne lui rapporte rien, et si les pertes sont reportées, les premiers 50\,000~\$ de P\&L net de l'année suivante ne rapportent rien non plus.
Le même trader dans une société salariée toucherait un salaire de base et un bonus discrétionnaire, avec une variance plus faible et, en général,
une espérance totale plus basse (le chapitre~13 évalue de telles offres).
\end{example}

\section{Tutoriel~: la longue queue}\label{tut:in:medium-frequency-proprietary-shops:tut}

\begin{tutorial}
\textbf{Objectif.} Ajuster la distribution des tailles des broker-dealers à partir des classes publiées par la FINRA et lire les entrées et les sorties.
\textbf{État final~:} la \cref{fig:in:medium-frequency-proprietary-shops:ccdf} et l'indice de queue avec son erreur type.
\begin{tutsteps}
\item \textbf{Les classes.} \texttt{data/industry/finra\_size\_bins.csv} contient les quatorze classes de taille de la FINRA pour 2020--2024~;
\texttt{in\_firmsize.bins(2024)} les lit sous forme d'enregistrements \texttt{firm.firmsize.Bin}.
\item \textbf{Ajustement.} \texttt{firm.firmsize.pareto\_fit(bins, 11)} maximise la vraisemblance sur classes
(le \cref{lst:in:medium-frequency-proprietary-shops:fit})~; \texttt{lognormal\_fit} ajuste la comparaison.
\omcode{code/firm/firmsize/firm_firmsize.py}{51}{58}{L'ajustement de Pareto sur classes et son erreur type à partir de
l'information observée.}\label{lst:in:medium-frequency-proprietary-shops:fit}
\item \textbf{Vérification.} \texttt{expected\_counts} et \texttt{chi2} comparent les deux ajustements~; il faut refaire l'ajustement avec $x_{\min}$ égal à 16, 21,
31 et 51, et voir $\alpha$ passer de 0.60 à 0.65.
\item \textbf{Flux.} \texttt{in\_firmsize.entry\_exit()} calcule les taux de sortie et d'entrée à l'origine de la
\cref{fig:in:medium-frequency-proprietary-shops:flows}.
\end{tutsteps}
\end{tutorial}

\textbf{Pour aller plus loin.} Ajuster les classes de 2020 et comparer l'indice avec celui de 2024 (l'exercice~7)~; simuler les tailles de sociétés à partir
de la queue ajustée avec \texttt{sample\_pareto} et vérifier que l'ajustement retrouve l'indice dont on est parti.

\section{Construction~: distributions de tailles à partir de comptages par classes}\label{bld:in:medium-frequency-proprietary-shops:firmsize}

\begin{build}
\bfield{Rôle} Ajuster des distributions de tailles aux comptages par classes publiés par les régulateurs, pour les tailles de sociétés ici et pour les bandes de rémunération
et les effectifs plus loin (chapitres~12 et~14).
\bfield{Interface} \texttt{firm.firmsize}~: \texttt{Bin(lo, hi, n)}~; \texttt{pareto\_fit(bins, x\_min)}~;
\texttt{lognormal\_fit(bins, x\_min)}~; \texttt{expected\_counts}~; \texttt{chi2}~; \texttt{ccdf}~; \texttt{sample\_pareto}.
\bfield{Règles} Les bornes entières sont continues aux demi-unités~; une classe ouverte en haut a une probabilité égale à la queue au-delà de sa
borne inférieure~; un ajustement indique sa borne inférieure et son erreur type.
\bfield{Tests d'acceptation} \texttt{code/firm/firmsize/tests/}~: l'ajustement retrouve un indice connu à partir de 20\,000 tailles simulées
à trois erreurs types près~; les effectifs attendus somment au total observé~; la part de queue diminue avec la taille~;
l'ajustement lognormal s'exécute.
\bfield{Pour aller plus loin} Choisir $x_{\min}$ en minimisant la distance entre la queue ajustée et la queue observée (la méthode
de Clauset--Shalizi--Newman)~; un intervalle bootstrap pour $\alpha$ sur les classes.
\end{build}

\begin{omsources}
\item FINRA, Industry Snapshot 2025~: tableaux 1.1.3, 2.1.9, 2.2.3 et 2.6.1--2.6.3.
\item AFM (Autorité néerlandaise des marchés financiers), Market participants subject to a licence obligation, sur
les salles de trading partagées.
\item CFTC v.\ Traders Global Group Inc.\ et al., D.N.J.\ No.~23-cv-11808~: plainte (29 août 2023)~; rapport et
recommandation du special master (13 mai 2025)~; résumé de l'ordonnance du tribunal par Willkie Farr \& Gallagher (26 mai 2025).
\item A.~Clauset, C.~R.~Shalizi et M.~E.~J.~Newman, «~Power-law distributions in empirical data~», \emph{SIAM Review}
51(4), 2009.
\end{omsources}

\section{Exercices}

\begin{exercise}[$\star$]\label{exo:in:medium-frequency-proprietary-shops:1}
D'après l'encadré daté, quelle part des broker-dealers comptait plus de 500 représentants enregistrés à la fin de 2024, et
quelle part de l'ensemble des représentants employaient-ils~?
\end{exercise}

\begin{exercise}[$\star$]\label{exo:in:medium-frequency-proprietary-shops:2}
Dans l'exemple, que touche le trader sur un P\&L brut de 500\,000~\$, et sur 150\,000~\$~?
\end{exercise}

\begin{exercise}[$\star$]\label{exo:in:medium-frequency-proprietary-shops:3}
Sous une queue de Pareto avec $\alpha=0.6$ à partir de $x_{\min}=10.5$, quelle est la probabilité qu'une société comptant au moins onze
représentants en ait plus de 1\,000~?
\end{exercise}

\begin{exercise}[$\star\star$]\label{exo:in:medium-frequency-proprietary-shops:4}
Calculer les taux moyens de sortie et d'entrée des broker-dealers sur 2011--2024, ainsi que le taux de contraction nette qu'ils impliquent.
\end{exercise}

\begin{exercise}[$\star\star$]\label{exo:in:medium-frequency-proprietary-shops:5}
Pourquoi un indice de queue inférieur à un est-il dit lourd~? Quelle en est la conséquence pour la part des plus grandes
sociétés lorsque la population croît~?
\end{exercise}

\begin{exercise}[$\star\star$]\label{exo:in:medium-frequency-proprietary-shops:6}
À partir de la \cref{fig:in:medium-frequency-proprietary-shops:ccdf}, où la droite de Pareto se situe-t-elle au-dessus des parts observées,
et que dit cela des plus grandes sociétés~?
\end{exercise}

\begin{exercise}[$\star\star\star$]\label{exo:in:medium-frequency-proprietary-shops:7}
\emph{Programmation.} Ajuster les classes de 2020 à partir de $x_{\min}=11$. L'indice diffère-t-il de celui de 2024 de plus de deux erreurs types~?
\end{exercise}

\begin{exercise}[$\star\star\star$]\label{exo:in:medium-frequency-proprietary-shops:8}
\emph{Trouver la faille.} «~La moitié des broker-dealers comptent dix salariés ou moins, donc un poste de trading typique se trouve dans une société de
dix personnes.~»
\end{exercise}

\section{Problème~: la longue queue}

\begin{problem}[{Problème du week-end --- la longue queue}]\label{pb:in:medium-frequency-proprietary-shops:1}
Un chercheur ayant reçu une offre d'une société de \omterm{def:in:medium-frequency-proprietary-shops:mft}{trading à moyenne fréquence} de douze personnes veut comprendre où se situent
de telles sociétés dans l'industrie et ce que signifie cet arrangement.

\medskip\noindent\textbf{Partie I --- La catégorie.}
\begin{enumerate}
\item Définir le \omterm{def:in:medium-frequency-proprietary-shops:mft}{trading à moyenne fréquence} et le situer entre le trading à haute fréquence et un fonds systématique.
\item Pourquoi ses coûts et son infrastructure diffèrent-ils de ceux d'un teneur de marché~?
\item Quelle unité la FINRA utilise-t-elle pour la taille des sociétés, et pourquoi sous-estime-t-elle une société pour compte propre~?
\item Combien de broker-dealers la FINRA a-t-elle classés dans ses segments de négociation pour compte propre et de market making, et quelle était leur taille~?
\item Quelle part des broker-dealers comptait dix représentants enregistrés ou moins, et cinquante ou moins~?
\end{enumerate}

\medskip\noindent\textbf{Partie II --- La queue.}
\begin{enumerate}[resume]
\item Écrire la probabilité d'une classe sous une queue de Pareto avec des bornes d'une demi-unité.
\item Donner l'indice ajusté à partir de onze représentants et son erreur type.
\item Comparer les ajustements de Pareto et de loi log-normale au moyen de leurs distances du chi-deux.
\item Comment l'indice évolue-t-il avec la borne inférieure~?
\item Que disent les parts des représentants selon la taille des sociétés que l'indice dit aussi~?
\end{enumerate}

\medskip\noindent\textbf{Partie III --- Entrée, sortie et arrangements.}
\begin{enumerate}[resume]
\item Donner les taux moyens de sortie et d'entrée sur 2011--2024, ainsi que les effectifs de 2024.
\item Définir une \omterm{def:in:medium-frequency-proprietary-shops:arcade}{salle de trading partagée} (trading arcade) et dire ce que le régulateur néerlandais exige de l'une d'elles.
\item Définir un \omterm{def:in:medium-frequency-proprietary-shops:funded}{programme de trader financé} (funded trader) et dire pourquoi il n'est pas un employeur.
\item Énoncer, de façon neutre, les allégations de l'affaire du trader financé, son issue et sa date.
\item Calculer la rémunération du trader dans l'exemple pour un P\&L brut de 800\,000~\$, 200\,000~\$ et 100\,000~\$.
\end{enumerate}

\medskip\noindent\textbf{Partie IV --- Le verdict.}
\begin{enumerate}[resume]
\item Énoncer le \emph{résultat nommé}~: l'indice de queue avec son erreur type et son intervalle de variation lorsque la borne inférieure se déplace, et
la part des représentants dans les petites sociétés.
\item Nommer les quatre façons dont un poste dans une petite société diffère d'un poste dans une grande.
\item Que doit demander le chercheur sur le capital, la rémunération et les contrôles~?
\item Quelle est à peu près la probabilité qu'un broker-dealer quitte le registre au cours d'une année donnée~?
\item En deux phrases, comment le chercheur doit-il peser l'offre~?
\end{enumerate}
\end{problem}

\section{Questions d'entretien}

\begin{interviewq}[$\star$ \iqroles{chercheur, trader}]\label{iq:in:medium-frequency-proprietary-shops:1}
Qu'est-ce qui change, en recherche comme en infrastructure, lorsque la durée de détention d'une stratégie passe de quelques secondes à une journée~?
\end{interviewq}

\begin{interviewq}[$\star$ \iqroles{trader}]\label{iq:in:medium-frequency-proprietary-shops:2}
On vous propose 50~\% de votre P\&L net sans salaire, ou un salaire de 150\,000~\$ et un bonus de 10~\% du P\&L net.
À partir de quel P\&L net les deux options sont-elles équivalentes~?
\end{interviewq}

\begin{interviewq}[$\star\star$ \iqroles{chercheur}]\label{iq:in:medium-frequency-proprietary-shops:3}
Vous ne disposez que d'effectifs regroupés par classes de taille des sociétés. Comment estimez-vous un indice de queue, et
comment savez-vous si une loi de puissance est la bonne forme~?
\end{interviewq}

\begin{interviewq}[$\star\star$ \iqroles{risque}]\label{iq:in:medium-frequency-proprietary-shops:4}
Plusieurs traders opèrent sous l'agrément d'une même société, chacun avec son propre capital. Quels risques le titulaire de l'agrément supporte-t-il~?
\end{interviewq}

\begin{interviewq}[$\star\star$ \iqroles{chercheur}]\label{iq:in:medium-frequency-proprietary-shops:5}
Une stratégie de \omterm{def:in:medium-frequency-proprietary-shops:mft}{trading à moyenne fréquence} rapporte 5~pb par transaction avant coûts, change chaque position une fois par jour
sur 20~millions de dollars de capital avec un taux de rotation de 100~\% par jour, et paie 2~pb de coûts. Quel est son rendement net
annuel sur le capital, et qu'est-ce qui limite sa montée en échelle~?
\end{interviewq}

\begin{interviewq}[$\star\star\star$ \iqroles{chercheur, risque}]\label{iq:in:medium-frequency-proprietary-shops:6}
Une société vend des évaluations à des candidats traders et promet des partages de profits à ceux qui les réussissent.
Comment reconnaîtriez-vous, à partir de ses chiffres, si son chiffre d'affaires provient du trading ou des frais~?
\end{interviewq}
